\documentclass[12pt,oneside]{book}

% ============================================================
% METADATOS DEL DOCUMENTO
% ============================================================
\newcommand{\universidad}{Universidad de Buenos Aires}
\newcommand{\facultad}{Facultad de Derecho}
\newcommand{\materia}{Derecho de la Integración}
\newcommand{\titulo}{Tratados marco, ALADI, Mercosur y grados de integración económica}
\newcommand{\autor}{Isaac Edgar Camacho Ocampo}
\newcommand{\anio}{2026}

% ============================================================
% PAQUETES
% ============================================================
\usepackage[utf8]{inputenc}
\usepackage[T1]{fontenc}
\usepackage[spanish,es-nodecimaldot]{babel}

\usepackage[
    top=2.5cm,
    bottom=2.5cm,
    left=3cm,
    right=2.5cm,
    headheight=15pt
]{geometry}

\usepackage{amsmath}
\usepackage{amssymb}
\usepackage{mathtools}

\usepackage{graphicx}
\usepackage{float}
\usepackage{caption}
\usepackage{subcaption}

\usepackage{booktabs}
\usepackage{array}
\usepackage{longtable}
\usepackage{tabularx}

\usepackage{enumitem}

\usepackage{xcolor}
\usepackage[most]{tcolorbox}

\usepackage{fancyhdr}
\usepackage{ragged2e}

\usepackage[
    colorlinks=true,
    linkcolor=blue!50!black,
    citecolor=green!40!black,
    urlcolor=blue!60!black,
    pdftitle={\titulo},
    pdfauthor={\autor},
    pdfsubject={Apuntes universitarios},
    bookmarks=true
]{hyperref}

% ============================================================
% RUTAS DE IMÁGENES
% ============================================================
\graphicspath{
    {./img/}
    {./graficos/}
    {./figuras/}
}

% ============================================================
% CONFIGURACIÓN GENERAL
% ============================================================
\setcounter{tocdepth}{2}

\setlength{\parindent}{0pt}
\setlength{\parskip}{0.5em}

\setlist[itemize]{
    topsep=0.3em,
    itemsep=0.2em
}

\setlist[enumerate]{
    topsep=0.3em,
    itemsep=0.2em
}

\clubpenalty=10000
\widowpenalty=10000

% ============================================================
% ENCABEZADOS Y PIES
% ============================================================
\pagestyle{fancy}
\fancyhf{}

\fancyhead[L]{\nouppercase{\leftmark}}
\fancyhead[R]{\materia}

\fancyfoot[C]{\thepage}

\renewcommand{\headrulewidth}{0.4pt}
\renewcommand{\footrulewidth}{0pt}

\fancypagestyle{plain}{
    \fancyhf{}
    \fancyfoot[C]{\thepage}
    \renewcommand{\headrulewidth}{0pt}
}

% ============================================================
% COMANDOS MATEMÁTICOS
% ============================================================
\newcommand{\abs}[1]{\left|#1\right|}
\newcommand{\norm}[1]{\left\lVert#1\right\rVert}
\newcommand{\vect}[1]{\mathbf{#1}}
\newcommand{\deriv}[2]{\frac{d#1}{d#2}}
\newcommand{\pderiv}[2]{\frac{\partial#1}{\partial#2}}

% ============================================================
% CAJAS ACADÉMICAS
% ============================================================
\newtcolorbox{definition}[1]{
    enhanced,
    breakable,
    title={Definición: #1},
    fonttitle=\bfseries,
    colback=blue!3,
    colframe=blue!50!black,
    boxrule=0.8pt,
    arc=2mm
}

\newtcolorbox{important}{
    enhanced,
    breakable,
    title={Importante},
    fonttitle=\bfseries,
    colback=yellow!8,
    colframe=orange!70!black,
    boxrule=0.8pt,
    arc=2mm
}

\newtcolorbox{example}[1]{
    enhanced,
    breakable,
    title={Ejemplo: #1},
    fonttitle=\bfseries,
    colback=green!4,
    colframe=green!40!black,
    boxrule=0.8pt,
    arc=2mm
}

\newtcolorbox{formula}[1]{
    enhanced,
    breakable,
    title={Fórmula / Regla: #1},
    fonttitle=\bfseries,
    colback=purple!4,
    colframe=purple!50!black,
    boxrule=0.8pt,
    arc=2mm
}

\newtcolorbox{exercise}[1]{
    enhanced,
    breakable,
    title={Ejercicio: #1},
    fonttitle=\bfseries,
    colback=gray!5,
    colframe=gray!60!black,
    boxrule=0.8pt,
    arc=2mm
}

\newtcolorbox{note}{
    enhanced,
    breakable,
    title={Nota},
    fonttitle=\bfseries,
    colback=white,
    colframe=gray!50,
    boxrule=0.6pt,
    arc=2mm
}

% Entorno simple para citas textuales del docente (no es una de las seis cajas
% semánticas: se usa un bloque de cita estándar para no sobrecargar el texto
% de cajas, dado el alto número de citas presentes en el apunte original).
\newenvironment{citadocente}
  {\begin{quote}\itshape\RaggedRight}
  {\end{quote}}

\begin{document}

% ============================================================
% PORTADA
% ============================================================
\begin{titlepage}

\thispagestyle{empty}

\begin{center}

\vspace*{1cm}

{\large \textsc{\universidad}}

\vspace{0.2cm}

{\large \textsc{\facultad}}

\vspace{3cm}

{\Huge \textbf{\materia}}

\vspace{1cm}

{\LARGE \titulo}

\vspace{0.8cm}

\textit{Comprender la arquitectura de los tratados es aprender a leer, detrás de cada norma, el compromiso recíproco que sostiene a los pueblos.}

\vspace{1.2cm}

\rule{0.70\textwidth}{0.5pt}

\vspace{0.5cm}

{\large Apuntes de estudio}

\vspace{0.5cm}

\rule{0.70\textwidth}{0.5pt}

\vfill

{\large \autor}

\vspace{0.3cm}

{\large \anio}

\end{center}

\vfill

\begin{flushleft}
\textit{Este es un modesto aporte a los estudiantes de ``Derecho de la Integración'' no pretende sustituir el material oficial de la materia.}
\end{flushleft}

\end{titlepage}

% ============================================================
% ÍNDICE
% ============================================================
\tableofcontents

% ============================================================
% PRESENTACIÓN DE LA CLASE
% ============================================================
\chapter*{Presentación de la clase}
\addcontentsline{toc}{chapter}{Presentación de la clase}
\markboth{Presentación de la clase}{Presentación de la clase}

Esta clase retoma la distinción entre modelos de integración (formal e informal) y su carácter pluridimensional, para luego analizar en profundidad el marco de la Asociación Latinoamericana de Integración (ALADI): su antecedente, la ALALC; la cláusula de la nación más favorecida; los acuerdos de alcance parcial; y la relación entre el Mercosur, la Comunidad Andina y la ALADI. Se desarrolla también la diferencia práctica entre un tratado de integración y un tratado internacional común (con los casos del gas boliviano y el Convenio de Montreal), y se explican los distintos grados de integración económica: zona franca, unión fronteriza y zona de libre comercio, cerrando con una referencia a los procesos de integración no formales, en particular en el continente africano.

\section*{Relevancia profesional}
\addcontentsline{toc}{section}{Relevancia profesional}

Comprender la arquitectura de los tratados de integración no es un ejercicio teórico: es la herramienta que permite a un abogado determinar, ante un caso concreto de comercio exterior, inversión, transporte o circulación de personas, si una controversia se rige por el derecho de la integración (con sus órganos y procedimientos propios) o por el derecho internacional público o privado común. Distinguir un acuerdo de alcance parcial de un tratado de integración, o identificar cuándo la cláusula de la nación más favorecida es aplicable, es exactamente el tipo de análisis que se le exige a un asesor legal de una empresa que opera en el Mercosur o en la ALADI. Además, entender por qué ciertos procesos de integración fracasan (como ocurre con buena parte de las experiencias africanas) o por qué otros se profundizan (como la Comunidad Andina) forma el criterio profesional necesario para asesorar sobre riesgo jurídico e institucional en negocios internacionales.

\section*{Recorrido temático}
\addcontentsline{toc}{section}{Recorrido temático}

La clase se organiza como un recorrido en escalones, del mismo modo en que se avanza en cualquier proceso de integración real: primero se repasan los fundamentos conceptuales, luego se estudia el tratado marco que hoy rige la integración latinoamericana, la ALADI, y a partir de allí se aprende a distinguir un tratado de integración de un tratado internacional común. Con esas herramientas se analizan dos casos concretos de áreas de integración: el Mercosur y la Comunidad Andina. Después se ordenan los distintos grados de integración económica, de menor a mayor compromiso, y se cierra con el caso de las integraciones no formales, tomando África como ejemplo de un proceso que, en la mayoría de los casos, no logró consolidarse.

Un ejemplo cotidiano que atraviesa toda la clase es el de la familia: se compara pertenecer a un área de integración con formar una pareja o una familia. Así como no se puede estar casado con dos personas al mismo tiempo y cumplir al 100\% con ambas, un estado no puede, en principio, pertenecer a dos áreas de integración con el mismo objetivo. Y así como en una familia el bienestar de un hermano repercute en los demás, en un área de integración la crisis económica de un socio repercute en el resto por la interdependencia que genera el mercado único.

% ============================================================
% CAPÍTULO 1
% ============================================================
\chapter{Fundamentos de la integración}

\section{Repaso: modelos de integración e integración pluridimensional}

\textit{Pregunta guía: ¿Por qué la integración nunca es puramente económica?}

Se retoma lo trabajado en una clase anterior sobre distintos modelos de integración, formales e informales, y se recuerda que la integración puede ser cultural, social, política o económica. Por más que se defina un objetivo principal para un área de integración, no puede afirmarse que ese objetivo agote todos los efectos del proceso, porque inevitablemente va a influir también en otros aspectos, como el jurídico o el político.

\begin{citadocente}
``Por eso lo mencionábamos como que era un proceso pluridimensional, que afectaba distintos aspectos o dimensiones de los estados, del derecho y demás.''
\end{citadocente}

\begin{definition}{Integración pluridimensional}
Proceso que, aunque se lo defina o encuadre principalmente en una dimensión (por ejemplo, la económica), necesariamente repercute también en otras dimensiones del estado y del derecho (política, social, cultural, jurídica).
\end{definition}

Se agrega que adoptar una misma forma de integración no significa que los procesos sean idénticos en distintos continentes, porque existen factores culturales, sociales, geográficos y climáticos que condicionan su desarrollo.

\section{Factores geográficos y económicos que condicionan la integración}

\textit{Pregunta guía: ¿Cómo afecta la geografía al costo de integrarse?}

Para ilustrar cómo la geografía condiciona la integración, se compara la construcción de rutas entre países europeos, donde las distancias son cortas, con la situación de América Latina, donde las distancias son mucho mayores.

\begin{example}{Infraestructura vial y distancias}
Se compara el costo de infraestructura vial entre países europeos cercanos (España-Francia, España-Italia) con las distancias que debe cubrir la infraestructura en América Latina, señalando que la inversión estatal necesaria para facilitar el tránsito de personas y mercaderías resulta mucho más compleja en el segundo caso.

También se menciona el caso del puente entre Argentina y Paraguay en la zona de la triple frontera: aunque el tramo argentino y el paraguayo están construidos, sin financiamiento económico no se puede incentivar el comercio que ese puente permitiría.
\end{example}

Se cierra esta introducción recordando la consigna de buscar distintos tratados marco (es decir, tratados que crean áreas de integración) para identificar qué tipo de integración establecen, y cuáles son los medios y las finalidades que persiguen esos tratados.

% ============================================================
% CAPÍTULO 2
% ============================================================
\chapter{El marco de la ALADI}

\section{El Tratado de Montevideo de 1980 y la creación de la ALADI}

\textit{Pregunta guía: ¿Qué diferencia a la ALADI de su antecesora, la ALALC?}

Se identifica un tratado del cono sur de 1980 que abarca prácticamente toda América Latina, orientado a liberar el comercio: la sigla ALADI corresponde a Asociación Latinoamericana de Integración.

\begin{important}
\textbf{Norma citada.} Tratado de Montevideo de 1980 — instrumento fundacional de la Asociación Latinoamericana de Integración (ALADI), aún vigente. No debe confundirse con los tratados de Montevideo de derecho internacional privado (de 1889 y 1940), que regulan materias distintas.
\end{important}

La ALADI, si bien incluye el libre comercio, no aspira solamente a constituir una zona de libre comercio, sino a lograr una integración más profunda: un desarrollo social, económico y cultural más global. El objetivo principal del tratado es establecer, en forma gradual y progresiva, un mercado común latinoamericano, y la finalidad última es el desarrollo económico y social para asegurar un mejor nivel de vida de los pueblos.

\section{Fin versus medio en los tratados marco}

\textit{Pregunta guía: ¿Cómo se distingue el objetivo final de los instrumentos para alcanzarlo?}

Se marca una distinción metodológica central: cuando un tratado habla de ``establecer en forma gradual y progresiva un mercado común latinoamericano'', el mercado común no es el fin del tratado, sino el medio para alcanzar un fin más amplio, que es la integración latinoamericana progresiva.

\begin{citadocente}
``A veces los legisladores legislan las palabras, que parece que no se quedan con el derecho\ldots si vos me decís `con el objetivo de' o `con la finalidad de', me estás poniendo una finalidad; si me decís `a través de la constitución de un mercado común', es un medio.''
\end{citadocente}

\begin{definition}{Distinción fin/medio}
El fin de un tratado marco de integración es el objetivo último que se busca (por ejemplo, el desarrollo o la integración latinoamericana progresiva). El medio es el instrumento jurídico o institucional elegido para alcanzarlo (por ejemplo, la constitución de un mercado común). Lo que no debería cambiar entre distintos tratados es la finalidad; lo que puede variar es el medio elegido.
\end{definition}

Esta distinción es la razón por la cual se pide identificar, en cada tratado, cuál es el fin último perseguido y cuál es el medio institucional elegido para lograrlo, ya que ambos elementos suelen aparecer mezclados en la redacción de los tratados.

\section{La cláusula de la nación más favorecida}

\textit{Pregunta guía: ¿Por qué un estado no puede beneficiar a un solo socio comercial?}

El Tratado de Asunción constituye el antecedente del Mercosur y también tiene como objeto la constitución de un mercado común. Retomando la distinción entre fin y medio: el fin sería la integración latinoamericana progresiva, mientras que el mercado común del Mercosur se constituye como un medio para lograr esa integración mayor, optimizando la parte productiva, las economías de escala y la complementación comercial.

A partir de allí, el Mercosur es, en rigor, una subregión dentro de un área de integración más amplia, que es la ALADI: la ALADI se creó justamente con la idea de impulsar la integración latinoamericana en su conjunto, y el Mercosur representa el impulso que dieron cuatro estados dentro de ese marco mayor.

\begin{example}{La analogía familiar de los regalos}
``Si fuiste de compra y le compraste zapatos solamente a uno solo de tus hijos, se va a matar, ¿entienden? Para todos es igual.'' Con esta imagen se explica que, dentro de un área de integración, si un estado le otorga beneficios específicos a otro estado miembro, está obligado a extenderlos a todos los demás miembros, salvo que exista un acuerdo de alcance parcial o un área de integración específicamente aprobada que lo autorice a limitar ese beneficio a determinados estados.
\end{example}

\begin{definition}{Cláusula de la nación más favorecida}
Dentro de un espacio de integración, un estado miembro no puede otorgar beneficios específicos a otro estado sin extenderlos al resto de los miembros, salvo que dichos beneficios estén amparados por un acuerdo de alcance parcial o un área de integración debidamente autorizada dentro del marco general.
\end{definition}

\section{Acuerdos de alcance parcial y acuerdos bilaterales}

\textit{Pregunta guía: ¿Qué controla la secretaría de la ALADI sobre estos acuerdos?}

La ALADI crea, entre sus órganos, una secretaría general, una de cuyas funciones es permitir y controlar acuerdos bilaterales o multilaterales entre los estados parte de la ALADI. Para explicar por qué un estado no puede pertenecer, en principio, a dos áreas de integración con el mismo objetivo, se recurre nuevamente a una analogía familiar.

\begin{example}{La analogía del doble matrimonio}
Pertenecer a dos áreas de integración se compara con casarse dos veces: ``¿A los hijos de qué matrimonio mandas a la escuela privada y a cuáles a la escuela estatal si no te alcanza la plata? ¿Con quién vas de vacaciones?''. Según la doctrina, un estado no puede pertenecer a dos áreas de integración con el mismo objetivo, salvo que una sea, por ejemplo, política y la otra social o cultural, es decir, de una naturaleza distinta.
\end{example}

La ALADI no permite hacer acuerdos de integración con estados ajenos al bloque (esos deben regirse por el derecho internacional común), pero sí permite a sus miembros celebrar acuerdos bilaterales o multilaterales entre ellos, sean de integración o de alcance parcial, siempre que no vayan en contra de los principios y objetivos generales de la ALADI.

El origen histórico del Mercosur se enmarca en este esquema: con el retorno de la democracia en la década de 1980, Argentina y Brasil firmaron numerosos acuerdos bilaterales de alcance parcial, limitados al ámbito comercial, que fueron el antecedente directo de la creación del Mercosur. Posteriormente, cuando se propuso transformar ese vínculo bilateral en un tratado de integración entre Argentina y Brasil, Uruguay y Paraguay solicitaron sumarse, dando origen al Tratado de Asunción entre los cuatro países, siempre dentro del ámbito que permite la ALADI.

Finalmente, la secretaría de la ALADI registra y controla los distintos acuerdos de alcance parcial y los acuerdos de integración existentes entre los estados miembros, verificando que no contraríen los objetivos generales del bloque, precisamente por la vigencia de la cláusula de la nación más favorecida.

\begin{table}[H]
\centering
\caption*{\textbf{Cuadro Cornell --- Capítulo 1 y 2}}
\begin{tabularx}{\textwidth}{
    >{\raggedright\arraybackslash}p{0.28\textwidth}
    >{\raggedright\arraybackslash}X
}
\toprule
\textbf{Preguntas / palabras clave} &
\textbf{Notas / respuestas} \\
\midrule

\textbf{¿Qué significa que la integración sea un proceso pluridimensional?} &
Que, aunque se defina un objetivo principal (por ejemplo, económico), el proceso de integración inevitablemente repercute también en otras dimensiones del estado y del derecho: política, social, cultural y jurídica. \\

\textbf{¿Qué es la ALADI y cuál es su antecedente?} &
La Asociación Latinoamericana de Integración, creada por el Tratado de Montevideo de 1980. Su antecedente es la ALALC (Asociación Latinoamericana de Libre Comercio), del Tratado de Montevideo de 1960, orientada exclusivamente al libre comercio. \\

\textbf{¿Cuál es la diferencia entre el fin y el medio en un tratado de integración?} &
El fin es el objetivo último perseguido (por ejemplo, la integración latinoamericana progresiva). El medio es el instrumento institucional elegido para lograrlo (por ejemplo, constituir un mercado común). El medio puede variar; la finalidad, en principio, no debería cambiar. \\

\textbf{¿En qué consiste la cláusula de la nación más favorecida?} &
Establece que un estado miembro de un área de integración no puede otorgar beneficios específicos a otro estado miembro sin extenderlos a los demás, salvo que existan acuerdos de alcance parcial o áreas de integración específicamente autorizadas. \\

\textbf{¿Qué controla la secretaría general de la ALADI?} &
Registra y controla los acuerdos de alcance parcial y los acuerdos de integración (bilaterales o multilaterales) que celebran los estados miembros entre sí, verificando que no contradigan los objetivos generales del bloque. \\

\midrule

\multicolumn{2}{p{\dimexpr\textwidth-2\tabcolsep\relax}}{
\textbf{Resumen:} La integración es un proceso pluridimensional que excede lo económico y repercute en lo político, social, jurídico y cultural. Dentro de ese marco, la ALADI —creada por el Tratado de Montevideo de 1980 como heredera de la ALALC de 1960— constituye hoy el marco general de la integración latinoamericana, sujeto a la cláusula de la nación más favorecida y al control de su secretaría general.
} \\

\bottomrule
\end{tabularx}
\end{table}

% ============================================================
% CAPÍTULO 3
% ============================================================
\chapter{Integración vs. derecho internacional común}

Distinguir un tratado de integración de un tratado internacional común no siempre es sencillo en la práctica, aunque en la teoría la diferencia parezca clara. Para trabajar esta distinción se analizan tres ejemplos concretos.

\section{El caso del gas boliviano}

\textit{Pregunta guía: ¿Por qué una compraventa internacional no es un tratado de integración?}

\begin{example}{Compraventa internacional de gas}
Durante el gobierno anterior, la Argentina no producía gas suficiente y le compraba gas a Bolivia. El tratado de compraventa de gas firmado con Bolivia es un contrato internacional de derecho privado, no un acuerdo enmarcado en un área de integración. Esto significa que Bolivia no tenía obligación de venderle gas a la Argentina al mismo precio al que se lo vendía a otro estado, porque el vínculo no estaba alcanzado por la cláusula de la nación más favorecida propia de un área de integración.
\end{example}

La clave para distinguir este caso de un acuerdo de integración está en la finalidad económica de cada parte: la Argentina buscaba proveerse de gas para su propio estado, y Bolivia buscaba cobrar por su producto. No existía una finalidad de beneficio recíproco propia de un área de integración, sino un intercambio de intereses individuales típico de un contrato internacional de compraventa.

\begin{citadocente}
``Yo buscaba mi interés y vos buscabas el tuyo\ldots no era que mi intención era, sabés qué, una mano lava la otra, las dos lavan la cara.''
\end{citadocente}

\section{El Convenio de Montreal y el transporte aéreo}

\textit{Pregunta guía: ¿Alcanza con regular un aspecto técnico para hablar de integración?}

Se plantea la consulta por el Convenio de Montreal, que regula las relaciones internacionales aéreas de las compañías en materia aeronáutica.

\begin{important}
\textbf{Norma citada.} Convenio de Montreal — instrumento internacional en materia aeronáutica que regula las relaciones internacionales aéreas de las compañías de transporte. No constituye un tratado de integración, sino una regulación técnica del tráfico aéreo internacional.
\end{important}

\begin{example}{El tránsito vehicular como analogía}
``¿Cómo modifica el tráfico aéreo en los espacios internacionales? Vamos a hacer un acuerdo para organizarnos un poco en ese ámbito.'' Esta regulación se compara con establecer que una calle es de mano única o contramano: se trata de ordenar un aspecto técnico específico (el tráfico aéreo), sin que ello implique una interdependencia económica entre los estados firmantes.
\end{example}

Frente a la objeción de que la finalidad declarada del convenio también podía leerse como una integración, se retoma el concepto de interdependencia: la integración supone que la actitud económica de un estado repercute necesariamente en la economía del otro, porque el mercado es único. En el caso del Convenio de Montreal, en cambio, lo único que se regula es que los aviones no colisionen entre sí; no existe un objetivo de que un estado incentive el área de vuelo del otro en el sentido de una interdependencia económica plena.

\begin{definition}{Interdependencia}
Elemento distintivo de un área de integración económica. Significa que la actitud económica que adopta un estado (por ejemplo, aumentar o restringir importaciones) repercute directamente en la economía de los demás estados del bloque, porque comparten un mercado único. Sin ese efecto recíproco, no puede hablarse propiamente de integración, aunque exista una regulación internacional común sobre la materia.
\end{definition}

Se cierra el punto señalando que si un acuerdo se estructura, por ejemplo, en torno al turismo entre dos estados —estableciendo libre circulación e igualdad de trato con fines turísticos— sí podría constituir una integración económica, aunque focalizada en un solo servicio (el turístico) en lugar de abarcar todos los bienes y servicios, porque en ese caso sí existe un beneficio económico recíproco entre ambos estados.

\section{El tratado de cooperación judicial del Mercosur}

\textit{Pregunta guía: ¿Qué rol cumplen las normas derivadas dentro de un área ya integrada?}

Se presenta un tratado orientado a la cooperación en materia civil, comercial, administrativa y laboral, vigente dentro del ámbito del Mercosur. Este instrumento no constituye un tratado marco (es decir, no es el que crea el área de integración), sino que se trata de una norma elaborada por los órganos del Mercosur ya existente, que regula un aspecto específico: la cooperación judicial internacional en materia civil, comercial, administrativa y laboral.

\begin{important}
\textbf{Norma citada.} Tratado sobre cooperación judicial internacional del Mercosur — norma de derecho derivado (elaborada por los órganos propios del bloque, no un tratado fundacional) que regula la cooperación en materia civil, comercial, administrativa y laboral entre los estados miembros.
\end{important}

\begin{definition}{Derecho derivado}
Conjunto de normas elaboradas por los órganos propios de un área de integración ya constituida (a diferencia del tratado marco, que es el instrumento fundacional). El tratado de cooperación judicial del Mercosur se califica como derecho derivado y no como un tratado fundacional.
\end{definition}

\begin{table}[H]
\centering
\caption*{\textbf{Cuadro Cornell --- Capítulo 3}}
\begin{tabularx}{\textwidth}{
    >{\raggedright\arraybackslash}p{0.28\textwidth}
    >{\raggedright\arraybackslash}X
}
\toprule
\textbf{Preguntas / palabras clave} &
\textbf{Notas / respuestas} \\
\midrule

\textbf{¿Por qué el contrato de compraventa de gas entre Argentina y Bolivia no es un tratado de integración?} &
Porque es un contrato internacional de derecho privado en el que cada parte persigue su propio interés (Argentina proveerse de gas, Bolivia cobrar por su producto), sin la interdependencia económica recíproca propia de un área de integración. \\

\textbf{¿Por qué el Convenio de Montreal no constituye, por sí solo, un tratado de integración?} &
Porque regula un aspecto técnico específico (el tráfico aéreo internacional) para evitar colisiones y ordenar el espacio aéreo, sin generar la interdependencia económica recíproca que caracteriza a un área de integración. \\

\textbf{¿Qué es la interdependencia en materia de integración económica?} &
El fenómeno por el cual la actitud económica de un estado repercute directamente en la economía de los demás estados del bloque, porque comparten un mercado único. \\

\midrule

\multicolumn{2}{p{\dimexpr\textwidth-2\tabcolsep\relax}}{
\textbf{Resumen:} Distinguir un tratado de integración de un tratado internacional común exige identificar si existe o no interdependencia económica recíproca entre los estados, más allá de que ambos tipos de instrumentos puedan tratar aspectos económicos.
} \\

\bottomrule
\end{tabularx}
\end{table}

% ============================================================
% CAPÍTULO 4
% ============================================================
\chapter{Mercosur y Comunidad Andina}

\section{El Tratado de Asunción dentro del ámbito de la ALADI}

\textit{Pregunta guía: ¿Por qué el Mercosur se define como subregión de la ALADI?}

El Mercosur se define a sí mismo, en su propio discurso fundacional, como una integración tendiente a una integración mayor dentro del ámbito de la ALADI. Es decir, Argentina, Brasil, Uruguay y Paraguay se agruparon para darle un impulso específico al proceso, pero dejando la puerta abierta a que otros estados de la región se sumaran, ya que la finalidad última sigue siendo la integración latinoamericana en su conjunto que persigue la ALADI.

Un acuerdo bilateral que fuera en contra de los objetivos del bloque no podría calificarse como acuerdo de integración, sino que quedaría reducido al ámbito puramente internacional, como un contrato de compraventa común (en el sentido del ejemplo del gas boliviano desarrollado en el capítulo anterior).

\section{Del Grupo Andino a la Comunidad Andina}

\textit{Pregunta guía: ¿Cómo evolucionó el bloque andino frente al peso de Argentina y Brasil?}

Se presenta el antecedente del Grupo Andino, tratado de 1969; en el programa de la materia todavía figura con ese nombre porque no había sido actualizado, aunque el bloque ya se denomina Comunidad Andina. El objetivo original del acuerdo era promover el desarrollo equilibrado y armónico de los países miembros mediante la integración y la cooperación económica y social, acelerar su crecimiento y la generación de empleo, y facilitar la participación de sus miembros en el proceso de integración regional, con miras a la formación gradual de un mercado común latinoamericano.

\begin{important}
\textbf{Norma citada.} Acuerdo de Cartagena (1969) — tratado fundacional del Grupo Andino (hoy Comunidad Andina), integrado originalmente por Bolivia, Colombia, Ecuador, Perú y Venezuela. Su objetivo declarado: promover el desarrollo equilibrado y armónico de los países miembros mediante la integración y la cooperación económica y social.
\end{important}

\begin{citadocente}
``Escuchaste el `desarrollo armónico y equilibrado'. Es decir, es mutuo, es recíproco. No es que a mí lo único que me importa es hacer el gas de mi estado y si vos te vas a la ruina, el problema es tuyo.''
\end{citadocente}

Durante la década de 1960, dentro de la ALALC (antecesora de la ALADI, orientada exclusivamente al libre comercio), los dos países con mayor peso económico en Latinoamérica eran Argentina y Brasil, que en los hechos imponían las reglas del juego. Los países andinos —Bolivia, Colombia, Ecuador, Perú y Venezuela— se agruparon para poder defenderse de esa influencia permanente, sin salir del proyecto de integración latinoamericana, dando origen al Grupo Andino.

A diferencia de la ALALC, que fracasó y tuvo que transformarse en la ALADI, el Grupo Andino fue evolucionando hasta convertirse en la actual Comunidad Andina, que hoy cuenta con órganos supranacionales propios, entre ellos un tribunal de justicia supranacional. Este desarrollo institucional resulta muy similar al que tuvo la Unión Europea, y los estados miembros de la Comunidad Andina sostienen y conocen activamente su normativa, a diferencia de lo que ocurre con el conocimiento que los alumnos argentinos tienen sobre las normas del Mercosur.

\begin{citadocente}
``Si yo les pregunto ahora a ustedes normas del Mercosur, salvo el convenio de las líneas que trajo la compañera, no tienen la más mínima idea de una sola norma del Mercosur\ldots la sociedad está comprometida con este proyecto, lo tienen más aceitado que nosotros.''
\end{citadocente}

\section{La interdependencia económica entre estados socios}

\textit{Pregunta guía: ¿Por qué la crisis de un socio afecta a todo el bloque?}

Se retoma el concepto de interdependencia para explicar por qué, dentro de un área de integración, resulta habitual que existan mecanismos de ayuda entre los estados miembros. Si existe libre tránsito de personas y uno de los estados atraviesa alta tasa de desempleo, se produce migración hacia los estados con mejores condiciones laborales, lo que a su vez puede generar desempleo en el estado receptor.

\begin{example}{El trabajador y la lógica del mercado único}
``Vos mirás una blusa en una tienda, la ves en la otra, comprás la que sale más barato. Es lógica pura.'' Del mismo modo, si en un estado socio hay libre tránsito de personas y su economía genera emigración de trabajadores dispuestos a aceptar peores condiciones, ello puede terminar afectando la economía y el empleo del estado receptor.
\end{example}

Esta es la razón por la cual, dentro de un área de integración, se sostiene que la crisis económica, la hiperinflación o el deterioro de un estado socio termina repercutiendo, tarde o temprano, en los demás miembros del bloque, precisamente por la interdependencia que genera el mercado único.

\section{Mecanismos de auxilio financiero: el caso de Grecia}

\textit{Pregunta guía: ¿Qué instrumentos usa un bloque para asistir a un socio en crisis?}

Para ilustrar cómo un área de integración crea oportunidades para asistir a un estado socio en dificultades, se menciona el ejemplo del Banco Europeo y del Banco de Inversiones de la Unión Europea, que otorgan préstamos a los estados miembros que lo necesitan, bajo determinadas condiciones de control.

\begin{example}{La crisis de deuda de Grecia}
La crisis de deuda de Grecia dentro de la Unión Europea obligó a los demás socios europeos a intervenir para asistirla financieramente, imponiéndole a la vez condiciones de control sobre el uso de esos fondos: ``Te prestamos plata para que la utilices en esto, para que salgas del hoyo, y vos te compraste más pala para cavar más profundo el hoyo. Yo te presto plata, pero como acreedor tuyo te voy a controlar.''
\end{example}

Este tipo de mecanismos de asistencia financiera entre socios normalmente no está previsto en el tratado fundacional del bloque, sino que forma parte del derecho derivado de la integración, es decir, de las normas e instituciones que los órganos del bloque van creando con posterioridad a la firma del tratado marco.

\begin{table}[H]
\centering
\caption*{\textbf{Cuadro Cornell --- Capítulo 4}}
\begin{tabularx}{\textwidth}{
    >{\raggedright\arraybackslash}p{0.28\textwidth}
    >{\raggedright\arraybackslash}X
}
\toprule
\textbf{Preguntas / palabras clave} &
\textbf{Notas / respuestas} \\
\midrule

\textbf{¿Cómo se relaciona el Mercosur con la ALADI?} &
El Mercosur es una subregión dentro del área de integración más amplia que constituye la ALADI. Se define a sí mismo como una integración tendiente a una integración latinoamericana mayor, dentro del ámbito de la ALADI. \\

\textbf{¿Cuál es el origen histórico del Grupo Andino (hoy Comunidad Andina)?} &
Surgió en 1969 cuando Bolivia, Colombia, Ecuador, Perú y Venezuela se agruparon para contrarrestar la influencia dominante de Argentina y Brasil dentro de la ALALC, sin abandonar el proyecto de integración latinoamericana. \\

\textbf{¿En qué se diferencia hoy la Comunidad Andina de otros bloques latinoamericanos?} &
Cuenta con órganos supranacionales propios, incluido un tribunal de justicia supranacional, con un desarrollo institucional comparable al de la Unión Europea. \\

\textbf{¿Qué rol cumplen mecanismos como el Banco de Inversiones europeo frente a la crisis de un socio (caso Grecia)?} &
Permiten otorgar asistencia financiera a un estado miembro en crisis, bajo condiciones de control, porque su deterioro económico repercutiría en el resto del bloque por la interdependencia del mercado único. Estos mecanismos suelen formar parte del derecho derivado, no del tratado fundacional. \\

\midrule

\multicolumn{2}{p{\dimexpr\textwidth-2\tabcolsep\relax}}{
\textbf{Resumen:} Dentro del marco de la ALADI conviven subregiones como el Mercosur —que se define a sí mismo como un paso hacia una integración latinoamericana mayor— y bloques más antiguos y consolidados como la Comunidad Andina, cuya interdependencia económica explica tanto los riesgos compartidos entre socios como los mecanismos de auxilio financiero creados para asistir a un estado en crisis.
} \\

\bottomrule
\end{tabularx}
\end{table}

% ============================================================
% CAPÍTULO 5
% ============================================================
\chapter{Grados de integración económica}

Se retoma la diferenciación de las distintas áreas de integración económica según el grado de compromiso que asumen los estados, comparando estos grados con escalones que deben transitarse progresivamente para alcanzar una integración más profunda.

\section{La zona franca como decisión unilateral}

\textit{Pregunta guía: ¿Por qué una zona franca no compromete a otro estado?}

Una zona franca no constituye, en sí misma, una integración, porque es una decisión unilateral de un solo estado, que facilita la instalación de personas e industrias mediante beneficios fiscales, con el fin de impulsar el desarrollo de áreas económicamente desfavorecidas dentro de su propio territorio.

\begin{definition}{Zona franca}
Decisión unilateral de un estado que otorga beneficios fiscales y facilidades para la instalación de personas e industrias en un área determinada de su territorio, con el objetivo de impulsar zonas desfavorecidas. No constituye una integración porque no compromete a ningún otro estado.
\end{definition}

\section{La unión fronteriza}

\textit{Pregunta guía: ¿Qué territorio abarca realmente una unión fronteriza?}

La unión fronteriza es el primer nivel o escalón dentro de los procesos de integración económica propiamente dichos. A diferencia de lo que su nombre podría sugerir, no implica la unión de los estados en su totalidad, sino que se limita a las áreas de frontera entre los países involucrados.

\begin{definition}{Unión fronteriza}
Primer nivel de integración económica. Establece políticas específicas de tránsito de personas, bienes o servicios limitadas exclusivamente a las áreas de frontera entre los estados involucrados (por ejemplo, determinadas provincias o regiones limítrofes), sin extenderse a la totalidad del territorio de cada estado.
\end{definition}

La delimitación concreta del área de frontera alcanzada por este tipo de acuerdo depende de la organización interna de cada estado: puede tratarse de una provincia entera (como las provincias de Formosa y Misiones en el caso argentino) o de una región más acotada, y en general responde a una decisión política consensuada con las provincias o regiones involucradas, ya que estas deben estar de acuerdo con la apertura del tránsito en su territorio.

\begin{example}{Unión fronteriza entre Chile y Argentina}
Se cita un tratado entre Chile y Argentina orientado a facilitar la integración regional de las provincias fronterizas. Una unión fronteriza no equivale a establecer una zona de libre comercio entre las capitales de ambos países (por ejemplo, Santiago de Chile y Buenos Aires), sino que responde a una necesidad de regular una realidad de hecho: la circulación cotidiana de personas y mercaderías que ya ocurre entre las poblaciones fronterizas.
\end{example}

\section{La zona de libre comercio y el arancel intrazona cero}

\textit{Pregunta guía: ¿Por qué la desgravación no puede ser igual para todos los productos?}

La zona de libre comercio se enfoca exclusivamente en el aspecto comercial —compra y venta de bienes— y, a diferencia de la unión fronteriza, abarca la totalidad del territorio de cada estado, no solo sus áreas limítrofes. Su aspiración es establecer aranceles de importación y exportación comunes, que se van reduciendo progresivamente hasta llegar al arancel intrazona cero.

\begin{definition}{Zona de libre comercio}
Grado de integración económica enfocado exclusivamente en el comercio de bienes, que abarca la totalidad del territorio de los estados miembros. Aspira a alcanzar el arancel intrazona cero, es decir, la eliminación progresiva y total de aranceles y restricciones al comercio entre los estados parte, sin beneficios fiscales especiales ni preferencias nacionales que favorezcan a terceros.
\end{definition}

\begin{citadocente}
``Lo que intento es llegar al arancel intrazona, dentro de la zona de integración, cero. Libre es libre de todo: libre de aranceles y libre de promoción.''
\end{citadocente}

La desgravación arancelaria no puede aplicarse de manera uniforme entre todos los estados ni sobre todos los productos por igual, porque la sensibilidad económica de cada rubro varía según el peso que tiene en el presupuesto de cada país.

\begin{example}{Tablas de desgravación diferenciadas}
Se compara el impacto de reducir un arancel sobre los cereales en la economía argentina frente a su impacto en la economía paraguaya, señalando que una misma rebaja del 5\% puede representar una pérdida presupuestaria de millones de dólares para un estado y ser insignificante para otro. Por eso, los órganos técnicos de un bloque de libre comercio suelen establecer tablas de desgravación diferenciadas: ``Te bajo un 5\% en los zapatos y vos me bajás un 3\%'', ajustando los porcentajes según la repercusión real en cada economía, y no de manera idéntica para todos los productos y estados.
\end{example}

Ningún estado puede eliminar de un día para el otro la totalidad de sus ingresos aduaneros, porque el estado se sostiene en gran medida gracias a los impuestos —incluida la educación y la salud pública—; por eso el proceso de desgravación es necesariamente gradual y está sujeto a permanente revisión por parte de los órganos técnicos del bloque.

\section{Complejidad institucional creciente}

\textit{Pregunta guía: ¿Qué estructura orgánica exige cada grado de integración?}

A medida que se avanza de la unión fronteriza hacia la zona de libre comercio, crece también la complejidad institucional y normativa necesaria para sostener el proceso.

\begin{definition}{Complejidad institucional creciente}
En una unión fronteriza puede bastar con la coordinación entre las oficinas de frontera e inmigración propias de cada estado, sin necesidad de un órgano específico y permanente. En cambio, una zona de libre comercio requiere una estructura orgánica integrada por representantes de todos los estados miembros, dedicada de manera permanente a revisar aranceles, presupuestos y tablas de desgravación, junto con una legislación específica en constante actualización.
\end{definition}

Se cierra este bloque recordando la distinción entre integración formal e informal: la integración formal se da cuando existe un tratado marco que instrumentaliza el proceso, crea los órganos correspondientes y define qué tipo de normas se aplican y cómo se generan sus efectos. Algunas áreas de integración, en cambio, se dan de hecho, sin ese tratado marco, especialmente en materia de defensa o de política exterior.

\begin{table}[H]
\centering
\caption*{\textbf{Cuadro Cornell --- Capítulo 5}}
\begin{tabularx}{\textwidth}{
    >{\raggedright\arraybackslash}p{0.28\textwidth}
    >{\raggedright\arraybackslash}X
}
\toprule
\textbf{Preguntas / palabras clave} &
\textbf{Notas / respuestas} \\
\midrule

\textbf{¿Qué es una zona franca y por qué no constituye integración?} &
Es una decisión unilateral de un estado que otorga beneficios fiscales para instalar personas e industrias en un área desfavorecida de su propio territorio. No compromete a ningún otro estado, por lo que no constituye un proceso de integración. \\

\textbf{¿Qué caracteriza a una unión fronteriza?} &
Es el primer nivel de integración económica. Establece políticas de circulación de personas, bienes o servicios limitadas exclusivamente a las áreas de frontera entre los estados, no a la totalidad de sus territorios. \\

\textbf{¿Qué es la zona de libre comercio y qué es el arancel intrazona cero?} &
Es un grado de integración enfocado en el comercio de bienes que abarca todo el territorio de los estados miembros. Aspira a eliminar progresivamente los aranceles entre los socios hasta llegar a un arancel intrazona cero, mediante tablas de desgravación diferenciadas según el impacto en cada economía. \\

\textbf{¿Por qué la desgravación arancelaria no puede ser idéntica para todos los productos y estados?} &
Porque el mismo porcentaje de rebaja arancelaria tiene un impacto presupuestario muy distinto según el peso relativo de cada producto en la economía de cada estado (por ejemplo, los cereales para Argentina frente a Paraguay). \\

\textbf{¿Cómo crece la complejidad institucional entre los distintos grados de integración?} &
A medida que se avanza de la unión fronteriza hacia la zona de libre comercio, se requiere una estructura orgánica más compleja: de la coordinación entre oficinas de frontera se pasa a órganos permanentes dedicados a revisar aranceles, presupuestos y normativa específica en constante actualización. \\

\midrule

\multicolumn{2}{p{\dimexpr\textwidth-2\tabcolsep\relax}}{
\textbf{Resumen:} Los grados de integración económica —zona franca, unión fronteriza y zona de libre comercio— constituyen escalones de compromiso creciente entre los estados, cada uno con mayor complejidad institucional y normativa.
} \\

\bottomrule
\end{tabularx}
\end{table}

% ============================================================
% CAPÍTULO 6
% ============================================================
\chapter{Integración no formal: el caso africano}

\section{Integración política y de defensa sin tratado marco}

\textit{Pregunta guía: ¿Puede haber integración sin un tratado formal?}

Existen procesos de integración que se producen de hecho, sin un tratado marco formal, particularmente en materia de defensa y de política exterior. Un ejemplo de ello es cuando un conjunto de estados vota sistemáticamente de manera conjunta en los organismos internacionales, lo cual revela una integración política o de defensa informal, aunque no exista un instrumento jurídico que la formalice.

\begin{definition}{Integración no formal}
Proceso de integración que se manifiesta en la práctica —por ejemplo, mediante el voto conjunto y sistemático de un grupo de estados en organismos internacionales, o mediante políticas de defensa compartidas— sin que exista un tratado marco que lo formalice institucionalmente. Este fenómeno es frecuente en el bloque africano.
\end{definition}

\section{Fracasos de integración en África: causas estructurales}

\textit{Pregunta guía: ¿Por qué la falta de un órgano supranacional no es la causa real del fracaso?}

Se dedica esta unidad a analizar por qué la gran mayoría de los procesos de integración en África no logran consolidarse en la práctica, a pesar de figurar formalmente creados.

\begin{citadocente}
``El noventa por ciento de ellas son un fracaso\ldots figuran de nombre, pero realmente no hay un efecto concreto.''
\end{citadocente}

Para explicar las causas de estos fracasos, se desarrolla el contexto histórico de los estados africanos: se trata, en su mayoría, de estados de formación reciente, que hasta hace pocas décadas eran colonias europeas (mencionando el caso de Francia e Inglaterra), y que dependieron durante mucho tiempo, económica, legislativa y políticamente, de las decisiones de las potencias coloniales, además de reunir en un mismo territorio culturas internamente muy distintas entre sí.

\begin{example}{La analogía del padre colonial}
``Porque era colonia, papá no te dejaba pelear, pero una vez que se fue papá, no bajamos la guardia.'' Con esta imagen se explica que la salida del dominio colonial dejó al descubierto tensiones internas entre culturas distintas que antes eran contenidas por la potencia colonizadora.
\end{example}

En cuanto a las causas económicas del fracaso, muchos estados africanos carecen de riqueza propia consolidada y presentan una economía con mucha precariedad. Un ejemplo es la producción ganadera: aunque un país africano cría ganado, si sus estándares fitosanitarios no cumplen los parámetros exigidos por los mercados europeo o americano, sus productos no logran acceder a esos mercados, lo cual limita drásticamente sus posibilidades de desarrollo comercial.

\begin{citadocente}
``Tu mercadería en mi estado no pasa los controles\ldots entonces no tengo industrias, no te puedo vender maquinarias. ¿Qué te puedo vender? Trabajo humano de la tierra, del campo, y nos lo vendemos entre nosotros.''
\end{citadocente}

La ausencia de una base económica sólida en los estados involucrados dificulta seriamente la consolidación de la integración, más allá de la voluntad política declarada en los tratados. Se cuestiona además una explicación frecuente en ámbitos académicos, según la cual el fracaso de estos procesos se debe a la falta de órganos supranacionales.

\begin{citadocente}
``Me he cansado de escuchar conferencias de profesores que dicen: es porque no tienen un órgano supranacional. ¿De qué estás hablando? Si no tienen para apagar la luz, no tienen para comprarse la casa. La propuesta que me estás haciendo para solucionar el problema es inviable, no tiene que ver con que el órgano sea supranacional.''
\end{citadocente}

Como contraejemplo dentro del propio continente africano, se menciona el caso de la Comunidad Económica y Monetaria de África Central (identificada en clase como ``la comunidad africana''), presentada como el proceso más avanzado de la región, cuyo Banco Central, sin embargo, se encuentra bajo fuerte influencia de Francia, lo cual condiciona la autonomía real de las decisiones económicas de sus estados miembros.

% TODO: contenido incompleto o ambiguo en las notas originales — el apunte no precisa el nombre completo ni la composición actual de la ``Comunidad Económica y Monetaria de África Central'' más allá de lo consignado en clase.

Se cierra la clase con la consigna de buscar, para el próximo encuentro, un tratado marco de integración —recomendando explorar tratados africanos y asiáticos, menos conocidos por los alumnos— con el fin de seguir practicando la identificación del grado de integración económica involucrado (zona de libre comercio, unión fronteriza, mercado común, unión económica, entre otros), remarcando que el objetivo del ejercicio es analizar y comprender el porqué de cada clasificación, y no memorizarla.

\begin{citadocente}
``A mí me gusta que analicemos, no que estudiemos de memoria\ldots lo que importa es que lo analicen, el porqué, se den cuenta el porqué, no lo repitamos como el trabajo práctico.''
\end{citadocente}

\begin{table}[H]
\centering
\caption*{\textbf{Cuadro Cornell --- Capítulo 6}}
\begin{tabularx}{\textwidth}{
    >{\raggedright\arraybackslash}p{0.28\textwidth}
    >{\raggedright\arraybackslash}X
}
\toprule
\textbf{Preguntas / palabras clave} &
\textbf{Notas / respuestas} \\
\midrule

\textbf{¿Qué es la integración no formal?} &
Un proceso de integración que se manifiesta en la práctica (por ejemplo, mediante el voto conjunto sistemático de un grupo de estados en organismos internacionales) sin que exista un tratado marco que lo formalice, frecuente en materia de defensa y política exterior. \\

\textbf{¿Cuáles son, según lo explicado en clase, las causas estructurales del fracaso de gran parte de los procesos de integración africanos?} &
La formación reciente de los estados, su prolongada dependencia colonial, la coexistencia de culturas internamente muy distintas, y sobre todo la precariedad económica de base, que impide sostener un compromiso recíproco real entre los socios, más allá de que exista o no un órgano supranacional. \\

\midrule

\multicolumn{2}{p{\dimexpr\textwidth-2\tabcolsep\relax}}{
\textbf{Resumen:} Existen procesos de integración no formales, sin tratado marco, especialmente en materia de defensa y política exterior; sin embargo, ningún tratado, por más órganos supranacionales que prevea, puede sostener una integración real si los estados involucrados carecen de una base económica sólida y de un compromiso recíproco genuino entre sus miembros, tal como evidencia el caso africano.
} \\

\bottomrule
\end{tabularx}
\end{table}

\end{document}
